\documentclass[10pt]{amsart}
\usepackage[a4paper,margin=23mm]{geometry}
\usepackage{amsmath,amssymb,amsthm,mathtools,hyperref,graphicx,xspace,letltxmacro,xpatch}
\usepackage[english]{babel}
\usepackage{xfrac,xcolor}

\usepackage{mathrsfs}
\let\mathcal\mathscr

\let\moins\smallsetminus
\DeclareMathOperator{\arccosh}{arccosh}
\DeclareMathOperator{\arctanh}{arctanh}
\DeclareMathOperator{\diam}{diam}
\DeclareMathOperator{\Fr}{Fr}
\DeclareMathOperator{\grad}{grad}
\DeclareMathOperator{\Homeo}{Homeo}
\DeclareMathOperator{\inte}{int}
\DeclareMathOperator{\Iso}{Iso}
\DeclareMathOperator{\Ric}{Ric}
\DeclareMathOperator{\SO}{SO}
\DeclareMathOperator{\rO}{O}
\DeclareMathOperator{\vol}{vol}
\newcommand{\II}{\mathrm{II}}
\newcommand{\pt}{\mathrm{pt}}
\let\parallel\Vert
\let\epsilon\varepsilon
\newcommand{\sep}{\;;\;}
\newcommand{\To}[1]{\mathchoice{\xrightarrow{\textstyle\kern4pt#1\kern3pt}}{\stackrel{#1}{\longrightarrow}}{}{}}
\newcommand\hto{\mathrel{\lhook\joinrel\to}}

\newcommand{\Psfrac}[2]{(\sfrac{#1}{#2})}
\newcommand{\psfrac}[2]{\sfrac{(#1)}{#2}}
\newcommand{\spfrac}[2]{\sfrac{#1}{(#2)}}
\newcommand{\Pspfrac}[2]{(\sfrac{#1}{(#2)})}

\DeclareFontFamily{U}{mathx}{\hyphenchar\font45}
\DeclareFontShape{U}{mathx}{m}{n}{
      <5> <6> <7> <8> <9> <10>
      <10.95> <12> <14.4> <17.28> <20.74> <24.88>
      mathx10
      }{}
\DeclareSymbolFont{mathx}{U}{mathx}{m}{n}
\DeclareFontSubstitution{U}{mathx}{m}{n}
\DeclareMathAccent{\wideparen}    {0}{mathx}{"75}

\newtheorem{theorem}{Theorem}[section]
\newtheorem{lemma}[theorem]{Lemma}
\newtheorem{proposition}[theorem]{Proposition}
\newtheorem{corollary}[theorem]{Corollary}

\theoremstyle{definition}
\newtheorem{definition}[theorem]{Definition}

\theoremstyle{remark}
\newtheorem{remark}[theorem]{Remark}
\newtheorem{example}[theorem]{Example}

\numberwithin{equation}{section}
\newtheorem{conjecture}[theorem]{Conjecture}

\newcommand{\RR}{\mathbb R}
\newcommand{\ZZ}{\mathbb Z}

\let\dpl\displaystyle
\providecommand{\skpt}{\par\smallskip\noindent}
\providecommand{\pointir}{.\ }
\providecommand{\eg}{e.g.}\providecommand{\ie}{i.e.}\providecommand{\etc}{etc.}\providecommand{\resp}{resp.}\providecommand{\cf}{cf.}
\providecommand{\mainmatter}{}\providecommand{\backmatter}{}
\let\Subsection\subsection\let\Subsubsection\subsubsection
\setlength{\emergencystretch}{2em}
